\documentclass[11pt,a4paper]{scrartcl} 

% Kalbos nustatymai (lietuvių kalbai)


% Įkeliame jūsų .sty failą (įsitikinkite, kad Overleaf projekte šalia yra failas manostilius.sty)
\usepackage{../manostilius_lt}

% Projekto informacija (būtina preambulėje dėl jūsų stiliaus specifikos)
\title{Laipsnis su natūraliuoju rodikliu ir jo savybės}
\author{Andrius Berniukevičius} 

\begin{document}

\maketitle


\section{Laipsnio apibrėžimas}

\begin{definition}
   Skaičiaus $a$ \vocab{laipsniu su natūraliuoju rodikliu $n$} (kai $n > 1$) vadinama $n$ vienodų dauginamųjų, kurių kiekvienas lygus $a$, sandauga.
\end{definition}

Užrašoma:
$$a^n = \underbrace{a \cdot a \cdot a \cdot \dots \cdot a}_{n \text{ kartų}}$$
Jei laipsnio rodiklis $n = 1$, tai skaičiaus pirmuoju laipsniu laikomas pats skaičius: $$a^1 = a$$
Užraše $a^n$:
\begin{itemize}
    \item \textbf{$a$} vadinamas \textbf{laipsnio pagrindu} (tai skaičius, kurį dauginame).
    \item \textbf{$n$} vadinamas \textbf{laipsnio rodikliu} (tai natūralusis skaičius, kuris rodo, kiek kartų pagrindas kartojasi sandaugoje).
\end{itemize}

\begin{remark}
    Svarbu nesupainioti daugybos ir kėlimo laipsniu:
\begin{itemize}
    \item \textbf{Daugyba:} $4 \cdot 3 = 4 + 4 + 4 = 12$ (skaičių sudedame $3$ kartus)
    \item \textbf{Laipsnis:} $4^3 = 4 \cdot 4 \cdot 4 = 64$ (skaičių sudauginame $3$ kartus)
\end{itemize}
\end{remark}

\section{Laipsnio su natūraliuoju rodikliu savybės}
Visų žemiau esančių savybių įrodymai remiasi pagrindiniu laipsnio su natūraliuoju rodikliu apibrėžimu: 
$$a^n = \underbrace{a \cdot a \cdot \dots \cdot a}_{n \text{ kartų}}$$ 
\begin{property}[Laipsnių su vienodais pagrindais daugyba]
Kai dauginami laipsniai su vienodais pagrindais, pagrindas paliekamas tas pats, o rodikliai sudedami:
$$a^m \cdot a^n = a^{m+n}$$
\end{property}
\begin{proof}
Užrašykime abu sandaugos narius kaip vienodų dauginamųjų sandaugas pagal apibrėžimą:
$$a^m \cdot a^n = (\underbrace{a \cdot a \cdot \dots \cdot a}_{m \text{ kartų}}) \cdot (\underbrace{a \cdot a \cdot \dots \cdot a}_{n \text{ kartų}})$$
Kadangi skliaustai daugyboje asociatyvumo dėsnio pagrindu esmės nekeičia, matome, kad skaičius $a$ iš viso bendroje sandaugoje pasikartoja $m + n$ kartų:
$$a^m \cdot a^n = \underbrace{a \cdot a \cdot \dots \cdot a}_{m+n \text{ kartų}} = a^{m+n}$$
\end{proof}

\begin{property}[Laipsnių su vienodais pagrindais dalyba]
Kai dalijami laipsniai su vienodais pagrindais, pagrindas paliekamas tas pats, o iš dalinio rodiklio atimamas daliklio rodiklis:
$$a^m : a^n = a^{m-n} \quad (\text{kai } a \neq 0 \text{ ir } m > n)$$
\end{property}
\begin{proof}
Užrašykime dalybą trupmena ir išskleiskime abu laipsnius sandaugomis:
$$\frac{a^m}{a^n} = \frac{\underbrace{a \cdot a \cdot \dots \cdot a}_{m \text{ kartų}}}{\underbrace{a \cdot a \cdot \dots \cdot a}_{n \text{ kartų}}}$$
Kadangi skaitiklyje yra $m$ dauginamųjų, o vardiklyje – $n$ dauginamųjų, ir žinome, kad $m > n$, galime suprastinti po $n$ dauginamųjų $a$ tiek skaitiklyje, tiek vardiklyje. Vardiklyje visi $a$ išsiprastina (lieka $1$), o skaitiklyje lieka $m - n$ dauginamųjų:
$$\frac{a^m}{a^n} = \underbrace{a \cdot a \cdot \dots \cdot a}_{m-n \text{ kartų}} = a^{m-n}$$
$\square$
\end{proof}
\begin{example}
$$\frac{a^7}{a^3} = \frac{a \cdot a \cdot a \cdot a \cdot a \cdot a \cdot a}{a \cdot a \cdot a}$$
Dabar pradedame prastinti vienodus dauginamuosius skaitiklyje ir vardiklyje poromis. Kadangi vardiklyje yra trys $a$, jie visi susiprastina su trimis $a$ iš skaitiklio:

$$\frac{\cancelto{1}{a} \cdot \cancelto{1}{a} \cdot \cancelto{1}{a} \cdot a \cdot a \cdot a \cdot a}{\cancelto{1}{a} \cdot \cancelto{1}{a} \cdot \cancelto{1}{a}} = \frac{1 \cdot 1 \cdot 1 \cdot a \cdot a \cdot a \cdot a}{1 \cdot 1 \cdot 1}$$
Vardiklyje lieka $1$ (jis pranyksta), o skaitiklyje lieka tik $7 - 3 = 4$ dauginamieji:

$$a \cdot a \cdot a \cdot a = a^4$$
\textbf{Išvada:} \\
Matome, kad iš pradinių 7 dauginamųjų skaitiklyje atėmėme (suprastinome) 3 dauginamuosius iš vardiklio, todėl galutinis laipsnio rodiklis yra gaunamas atimties būdu: $a^{7-3} = a^4$.
\end{example}    

\begin{property}[Laipsnio kėlimas laipsniu]
Keliant laipsnį laipsniu, pagrindas paliekamas tas pats, o rodikliai sudauginami:
$$(a^m)^n = a^{m \cdot n}$$
\end{property}
\begin{proof}
Pagal apibrėžimą, pagrindą $(a^m)$ reikia sudauginti iš savęs $n$ kartų:
$$(a^m)^n = \underbrace{a^m \cdot a^m \cdot \dots \cdot a^m}_{n \text{ kartų}}$$
Pritaikome pirmąją savybę (laipsnių daugybą) – pagrindas lieka $a$, o rodikliai sudedami:
$$(a^m)^n = a^{\underbrace{m + m + \dots + m}_{n \text{ kartų}}}$$
Kadangi skaičių $m$ sudedame iš savęs $n$ kartų, tai pagal daugybos apibrėžimą ši suma yra lygi $m \cdot n$:
$$(a^m)^n = a^{m \cdot n}$$
\end{proof}



\begin{property}[Sandaugos kėlimas laipsniu]
Keliant sandaugą laipsniu, kiekvienas dauginamasis keliamas tuo laipsniu atskirai:
$$(a \cdot b)^n = a^n \cdot b^n$$
\end{property}
\begin{proof}
Pagal apibrėžimą, sandaugą $(a \cdot b)$ sudauginame iš savęs $n$ kartų:
$$(a \cdot b)^n = \underbrace{(a \cdot b) \cdot (a \cdot b) \cdot \dots \cdot (a \cdot b)}_{n \text{ kartų}}$$
Remiantis daugybos perstatomumo (komutatyvumo) ir jungiamumo (asociatyvumo) dėsniais, sugrupuojame visus dauginamuosius $a$ ir visus $b$ atskirai:
$$(a \cdot b)^n = (\underbrace{a \cdot a \cdot \dots \cdot a}_{n \text{ kartų}}) \cdot (\underbrace{b \cdot b \cdot \dots \cdot b}_{n \text{ kartų}})$$
Pritaikę laipsnio apibrėžimą abiems sandaugoms skliaustuose, gauname:
$$(a \cdot b)^n = a^n \cdot b^n$$
\end{proof}
\begin{property}[Trupmenos kėlimas laipsniu]
Keliant trupmeną laipsniu, skaitiklis ir vardiklis keliami tuo laipsniu atskirai:
$$\left(\frac{a}{b}\right)^n = \frac{a^n}{b^n} \quad (\text{kai } b \neq 0)$$
\end{property}
\begin{proof}
Pagal apibrėžimą, trupmeną dauginame iš savęs $n$ kartų:
$$\left(\frac{a}{b}\right)^n = \underbrace{\frac{a}{b} \cdot \frac{a}{b} \cdot \dots \frac{a}{b}}_{n \text{ kartų}}$$
Pagal trupmenų daugybos taisyklę (skaitikliai sudauginami su skaitikliais, vardikliai su vardikliais):
$$\left(\frac{a}{b}\right)^n = \frac{\underbrace{a \cdot a \cdot \dots \cdot a}_{n \text{ kartų}}}{\underbrace{b \cdot b \cdot \dots \cdot b}_{n \text{ kartų}}}$$
Pritaikę laipsnio apibrėžimą skaitikliui ir vardikliui atskirai, turime:
$$\left(\frac{a}{b}\right)^n = \frac{a^n}{b^n}$$
\end{proof}

\section{Uždavinių sprendimo pavyzdžiai}

\begin{example}
Supaprastinkite reiškinį:
$$ \left(x^{45} : x^{25}\right)^2 \cdot\left(x^{5}\right)^3 : x^{10} $$
\textbf{Sprendimas:}
Atliekame veiksmus eilės tvarka, remdamiesi laipsnių savybėmis:
\begin{enumerate}
    \item Pirmiausia atliekame veiksmą skliaustuose (laipsnių dalyba):
    $$ x^{45} : x^{25} = x^{45-25} = x^{20} $$
    
    \item Gautą rezultatą keliame laipsniu (laipsnio kėlimas laipsniu):
    $$ (x^{20})^2 = x^{20 \cdot 2} = x^{40} $$
    
    \item Sutvarkome antrąjį narį (laipsnio kėlimas laipsniu):
    $$ (x^5)^3 = x^{5 \cdot 3} = x^{15} $$
    
    \item Atliekame daugybą (laipsnių daugyba):
    $$ x^{40} \cdot x^{15} = x^{40+15} = x^{55} $$
    
    \item Galiausiai atliekame dalybą:
    $$ x^{55} : x^{10} = x^{55-10} = x^{45} $$
\end{enumerate}

\textbf{Trumpas užrašas:}
$$ \left(x^{45} : x^{25}\right)^2 \cdot\left(x^{5}\right)^3 : x^{10} = (x^{20})^2 \cdot x^{15} : x^{10} = x^{40} \cdot x^{15} : x^{10} = x^{55} : x^{10} = x^{45} $$

\textbf{Atsakymas:} $x^{45}$.
\end{example}

\begin{example}
Supaprastinkite reiškinį:
$$ \frac{\left(\frac{243^5 \cdot 9^8}{81^3 \cdot 27^6}\right)^4 \cdot\left(3^{10}\right)^2}{\left(\frac{729^2}{9^4}\right)^3 \cdot 3^{12}} $$
\textbf{Sprendimas:} Atliekame veiksmus etapais, konvertuodami viskas į 3 laipsnius:

\begin{enumerate}
    \item Konvertuojame skaičius į 3 laipsnius ($243 = 3^5$, $9 = 3^2$, $81 = 3^4$, $27 = 3^3$):
    $$ \left(\frac{(3^5)^5 \cdot (3^2)^8}{(3^4)^3 \cdot (3^3)^6}\right)^4 = \left(\frac{3^{25} \cdot 3^{16}}{3^{12} \cdot 3^{18}}\right)^4 = \left(\frac{3^{41}}{3^{30}}\right)^4 = (3^{11})^4 = 3^{44} $$
    
    \item Skaičiuojame $(3^{10})^2 = 3^{20}$, todėl skaitiklis yra:
    $$ 3^{44} \cdot 3^{20} = 3^{64} $$
    
    \item Vardiklyje konvertuojame ($729 = 3^6$, $9 = 3^2$):
    $$ \left(\frac{729^2}{9^4}\right)^3 = \left(\frac{(3^6)^2}{(3^2)^4}\right)^3 = \left(\frac{3^{12}}{3^{8}}\right)^3 = (3^4)^3 = 3^{12} $$
    
    \item Visas vardiklis:
    $$ 3^{12} \cdot 3^{12} = 3^{24} $$
    
    \item Paskutinis žingsnis -- skaitiklis dalijamas iš vardiklio:
    $$ \frac{3^{64}}{3^{24}} = 3^{40} $$
\end{enumerate}

\textbf{Atsakymas:} $3^{40}$.
\end{example}

\section{Uždaviniai}
\begin{problem}
Atlikite veiksmus su laipsniais:

\begin{tasks}[label={(\arabic*)},label-offset=0.9em](2)
    \task $(((x^2)^5)^4)^3 : x^{100}$
    \task $\left(\frac{a^{25}}{b^{15}}\right)^4 \cdot \left(\frac{b^{30}}{a^{40}}\right)^2$
    \task $(y^{15} \cdot y^{25})^3 : (y^{12})^8 \cdot y^{10}$
    \task $(a^{10} \cdot b^{20} \cdot c^{30})^5 : (a^{12} \cdot b^{24} \cdot c^{36})^4$
    \task $\frac{(p^{18} : p^{10})^5 \cdot p^{20}}{(p^3 \cdot p^7)^4}$
    \task $\frac{(2a^{15})^4 \cdot (5a^{10})^2}{(10a^{30})^2}$
    \task $\left( \frac{x^{25} \cdot y^{35}}{x^{15} \cdot y^{20}} \right)^4$
    \task $\frac{u^{50}}{v^{30}} : \left(\frac{u^{15}}{v^{10}}\right)^3$
    \task $x^{100} : \left( \frac{(x^{15})^4 \cdot x^{20}}{x^{10}} \right)$
    \task $\frac{(m^{12} \cdot n^{18})^5 \cdot k^{40}}{(m^{20} \cdot n^{30})^3 \cdot (k^8)^5}$
    \task $(((z^3)^4)^2 \cdot z^{15}) : z^{30}$
    \task $\left(\frac{c^{40}}{d^{20}}\right)^3 : \left(\frac{c^{50}}{d^{30}}\right)^2$
    \task $\frac{(3x^{12}y^8)^4}{(9x^{20}y^{15})^2}$
    \task $\frac{(a^{15} \cdot a^{25})^4}{(a^{20})^7 \cdot a^{10}}$
    \task $(w^{50} : w^{20})^3 \cdot (w^{10} : w^5)^4$
    \task $\left( \frac{x^{30} \cdot y^{40} \cdot z^{50}}{x^{20} \cdot y^{25} \cdot z^{45}} \right)^5$
    \task $\frac{(2^{30} \cdot 2^{40})^2}{(2^{35})^4}$
    \task $(a^{18} \cdot b^{25})^5 : (a^{15} \cdot b^{20})^6$
    \task $\frac{(((p^5)^4)^3)^2}{p^{100} \cdot p^{15}}$
    \task $\left(\frac{m^{60}}{n^{40}}\right)^2 : \frac{m^{100}}{(n^{16})^5}$
    \task $\frac{(4^5 \cdot 8^3)^2}{16^6 \cdot 2^{10}}$
    \task $\frac{(9^{10} \cdot 27^5)^2}{(81^4)^3 \cdot 3^{15}}$
    \task $\frac{125^8 \cdot 25^{10}}{(5^6)^4 \cdot 25^5}$
    \task $\left( \frac{32^4 \cdot 8^6}{4^{10} \cdot 2^{15}} \right)^3$
    \task $\frac{(1000^5 \cdot 100^8)^2}{10^{50} \cdot 10000^2}$
    \task $\frac{64^5 \cdot (16^3)^4}{(4^8)^3 \cdot 8^5}$
\end{tasks}
\end{problem}

\newpage
\section{Atsakymai}
\begin{tasks}[label={(\arabic*)}, label-offset=0.9em](2)
    \task $x^{20}$
    \task $a^{20}$
    \task $y^{34}$
    \task $a^2 b^4 c^6$
    \task $p^{20}$
    \task $4a^{20}$
    \task $x^{40} y^{60}$
    \task $u^5$
    \task $x^{30}$
    \task $1$
    \task $z^9$
    \task $c^{20}$
    \task $x^8 y^2$
    \task $a^{10}$
    \task $w^{110}$
    \task $x^{50} y^{75} z^{25}$
    \task $1$
    \task $b^5$
    \task $p^5$
    \task $m^{20}$
    \task $16$
    \task $2187$
    \task $5^{10}$
    \task $512$
    \task $10000$
    \task $32768$
\end{tasks}




\end{document}